\section{Extension full subcategories}\label{s1}

\subsection{Extensions in Abelian categories}%\label{s1.1}

Let $\mathcal{A}$ be an Abelian category and $M,N\in\mathcal{A}$.
Recall that, for $d\in\mathbb{Z}_{\geq 0}$, the set $\mathrm{Ext}_{\mathcal{A}}^d(M,N)$ of {\em degree~$d$ extensions}
from~$M$ to~$N$ is def\/ined as the set of equivalence classes of exact sequences
\begin{gather*}
\xymatrix{
0\ar[r]&N\ar[r]&X_1\ar[r]&X_2\ar[r]&\cdots\ar[r]&X_d\ar[r]&M\ar[r]&0,
}
\end{gather*}
where all objects and all morphisms are in $\mathcal{A}$, modulo the minimal equivalence relation which contains the
binary relation given by existence of a~commutative diagram
\begin{gather*}
\xymatrix{
0\ar[r]&N\ar[r]\ar@{=}[d]&X_1\ar[r]\ar[d]&X_2\ar[r]\ar[d]&\cdots\ar[r]&X_d\ar[r]\ar[d]&M\ar[r]\ar@{=}[d]&0\\
0\ar[r]&N\ar[r]&Y_1\ar[r]&Y_2\ar[r]&\cdots\ar[r]&Y_d\ar[r]&M\ar[r]&0
}
\end{gather*}
The set $\mathrm{Ext}_{\mathcal{A}}^d(M,N)$ has the natural structure of an Abelian group via the Baer sum.
If $\mathcal{A}$ is $\Bbbk$-linear for some f\/ield $\Bbbk$, then $\mathrm{Ext}_{\mathcal{A}}^d(M,N)$ has the structure of
a~$\Bbbk$-vector space.
We refer to~\cite[Section~3.4]{We} for further information and details.

For $M\in\mathcal{A}$, the {\em projective dimension} $\mathrm{proj.dim}(M)$ of~$M$ is def\/ined as the maximal~$d$ such
that there is $N\in \mathcal{A}$ with $\mathrm{Ext}_{\mathcal{A}}^d(M,N)\neq 0$.
If such maximal~$d$ does not exist, then $\mathrm{proj.dim}(M):=\infty$.
Dually one def\/ines the {\em injective dimension} $\mathrm{inj.dim}(N)$ for $N\in\mathcal{A}$.
The {\em global dimension} $\mathrm{gl.dim}(\mathcal{A})\in\mathbb{Z}_{\geq 0}\cup\{\infty\}$ is the supremum of
projective dimensions taken over all objects in $\mathcal{A}$.
The global dimension coincides with the supremum of injective dimensions taken over all objects in $\mathcal{A}$
(see~\cite[Lemma~5.11.11]{Zi}).

\subsection{Extension full subcategories}\label{s1.2}

Let $\mathcal{A}$ be an Abelian category and $\mathcal{B}$ a~full Abelian subcategory of $\mathcal{A}$ in the sense that
the Abelian structure of $\mathcal{B}$ is inherited from $\mathcal{A}$.
In particular, the natural inclusion functor $\boldsymbol{\iota}:\mathcal{B}\to \mathcal{A}$ is exact.
Then, for every $M,N\in\mathcal{B}$ and every $d\in\mathbb{Z}_{\geq0}$, the functor $\boldsymbol{\iota}$ induces
homomorphisms
\begin{gather*}
\iota^d_{M,N}: \ \mathrm{Ext}_{\mathcal{B}}^d(M,N)\to \mathrm{Ext}_{\mathcal{A}}^d(M,N)
\end{gather*}
of Abelian groups.
In general, these homomorphisms $\iota^d_{M,N}$ are neither injective nor surjective.

We say that $\mathcal{B}$ is {\em extension full} in $\mathcal{A}$ provided that $\iota^d_{M,N}$ are bijective for all
$d\in\mathbb{Z}_{\geq0}$ and for all $M,N\in\mathcal{B}$.
We refer to~\cite[Section~2]{CoMa} for details.

Let $0\to K\to M\to N\to 0$ be a~short exact sequence in $\mathcal{B}$.
Then, for $Q\in \mathcal{B}$, application of $\mathrm{Hom}_{\mathcal{B}}(Q,{}_-)$ and
$\mathrm{Hom}_{\mathcal{A}}(Q,{}_-)$ to this short exact sequence produces the usual long exact sequences in homology
for the categories $\mathcal{B}$ and $\mathcal{A}$, respectively.
Moreover, the homomor\-phisms~$\iota^d_{Q,{}_-}$ give rise to a~homomorphism between these long exact sequences.
A~similar statement is true for $\mathrm{Hom}_{\mathcal{B}}({}_-,Q)$ and $\mathrm{Hom}_{\mathcal{A}}({}_-,Q)$.

\subsection{Checking extension fullness}%\label{s1.3}

In this section we formulate and prove three propositions which will be useful for our study of extension fullness later
in the paper.
The following statement is a~modif\/ication of~\cite[Lemma~4]{CoMa} which also allows for a~somewhat stronger formulation.

\begin{proposition}
\label{prop1}
Let $\mathcal{A}$ be an Abelian category and $\mathcal{B}$ a~Serre subcategory of $\mathcal{A}$.
Assume that $\mathcal{B}$ contains a~full subcategory $\mathcal{B}_0$ with the following properties:
\begin{enumerate}[$($a$)$]\itemsep=0pt
\item
\label{prop1.1}
Every object in $\mathcal{B}$ is a~quotient of an object in $\mathcal{B}_0$.
\item
\label{prop1.2}
The homomorphisms $\iota^d_{M,N}$ are bijective for all $d\in\mathbb{Z}_{\geq0}$, $M\in\mathcal{B}_0$ and
$N\in\mathcal{B}$.
\end{enumerate}
Then $\mathcal{B}$ is {\em extension full} in $\mathcal{A}$.
\end{proposition}

\begin{proof}
Our proof is similar to that of~\cite[Lemma~4]{CoMa}.
We prove the statement by induction on~$d$.
Since $\mathcal{B}$ is assumed to be a~Serre subcategory of $\mathcal{A}$, it is clear that $\iota^0_{Q,N}$ and
$\iota^1_{Q,N}$ are isomorphisms for all $Q,N\in\mathcal{B}$.

To prove the induction step, for $Q\in\mathcal{B}$ consider a~short exact sequence
\begin{gather*}
0\to K\to M\to Q\to 0
\end{gather*}
with $M\in\mathcal{B}_0$ and $K\in\mathcal{B}$, which exists by condition~\eqref{prop1.1}.
Applying $\mathrm{Hom}_{\mathcal{B}}({}_-,N)$ and $\mathrm{Hom}_{\mathcal{A}}({}_-,N)$ gives for each~$d$ the following
commutative diagram with exact rows:
\begin{gather*}
\xymatrix{
\mathrm{Ext}^{d-1}_{\mathcal{B}}(M,N)\ar[r]\ar[d]_{\iota^{d-1}_{M,N}}
&\mathrm{Ext}^{d-1}_{\mathcal{B}}(K,N)\ar[r]\ar[d]_{\iota^{d-1}_{K,N}}
&\mathrm{Ext}^d_{\mathcal{B}}(Q,N)\ar[r]\ar[d]_{\iota^d_{Q,N}}
& \mathrm{Ext}^d_{\mathcal{B}}(M,N)\ar[d]_{\iota^d_{M,N}}\\
\mathrm{Ext}^{d-1}_{\mathcal{A}}(M,N)\ar[r]&\mathrm{Ext}^{d-1}_{\mathcal{A}}(K,N)\ar[r]
&\mathrm{Ext}^d_{\mathcal{A}}(Q,N)\ar[r]& \mathrm{Ext}^d_{\mathcal{A}}(M,N)
}
\end{gather*}
Now, $\iota^{d-1}_{M,N}$ and $\iota^d_{M,N}$ are isomorphisms by condition~\eqref{prop1.2} and from the induction step
we have that $\iota^{d-1}_{K,N}$ is a~monomorphism.
The injective four lemma (by which we mean the statement of~\cite[Lemma~I.3.3(i)]{Mc}) therefore implies that
$\iota^d_{Q,N}$ is a~monomorphism.

It remains to show that $\iota^d_{Q,N}$ is an epimorphism.
For this we consider the following commutative diagram with exact rows:
\begin{gather*}
\xymatrix{
\mathrm{Ext}^{d-1}_{\mathcal{B}}(K,N)\ar[r]\ar[d]_{\iota^{d-1}_{K,N}}
&\mathrm{Ext}^{d}_{\mathcal{B}}(Q,N)\ar[r]\ar[d]_{\iota^{d}_{Q,N}}
&\mathrm{Ext}^d_{\mathcal{B}}(M,N)\ar[r]\ar[d]_{\iota^d_{M,N}}
& \mathrm{Ext}^d_{\mathcal{B}}(K,N)\ar[d]_{\iota^d_{K,N}}\\
\mathrm{Ext}^{d-1}_{\mathcal{A}}(K,N)\ar[r]&\mathrm{Ext}^{d}_{\mathcal{A}}(Q,N)\ar[r]
&\mathrm{Ext}^d_{\mathcal{A}}(M,N)\ar[r]& \mathrm{Ext}^d_{\mathcal{A}}(K,N)
}
\end{gather*}
As $\iota^{d-1}_{K,N}$ is a~bijection by the induction step, $\iota^d_{M,N}$ is a~bijection by condition~\eqref{prop1.2}
and $\iota^d_{K,N}$ is a~monomorphism by the previous paragraph, the surjective four lemma (by which we mean the
statement of~\cite[Lemma~I.3.3(ii)]{Mc}) implies that $\iota^d_{Q,N}$ is an epimorphism.
This completes the proof.
\end{proof}
